
\begin{table*}[h!]
    \centering
    \caption{Error taxonomy modified from~\citet{li-etal-2024-evaluating-mathematical}. Corresponding examples are in Fig.\ref{fig:error_type_examples} of appendix~\ref{sec:case_study}. Some infrequent error types are omitted to save space. For the full taxonomy, please check Tab.\ref{tab:error_taxonomy_full} in the appendix.}
    % \vspace{-2mm}
    \resizebox{\textwidth}{!}{%
    \begin{tabular}{l|l}
        \toprule[0.15em]
        \textbf{Error Type} & \textbf{Definition} \\
        \midrule
        \textbf{Calculation Error (CA)} & Error appears during the calculation process. \\
        \textbf{Counting Error (CO)} & Error occurs during the counting process. \\
        \textbf{Formula Confusion Error (FC)} & Error appears when applying formula in inappropriate scenario. \\
        \textbf{Question Misinterpretation Error(QM)} & Error appears because the \textbf{question} is misunderstood, such as ignoring specific constraints in the question. \\
        \textbf{Missing Step (MS)} & Error entails an incomplete generation of reasoning process, lacking a necessary step. \\
        \textbf{Confusing Concept Error (CC)} & Error occurs because two similar but actually different concepts are mistakenly confused. \\
        \textbf{Nonsensical Output (NO)} & Inconceivable, illogical, or question-irrelevant output.\\
        \bottomrule[0.15em]
    \end{tabular}%
    }
    \label{tab:error_taxonomy}
\end{table*}
